\chapter{指派、整数决策与全局界}
\label{ch:discrete}
\chaptertakeaway{区分可以通过线性规划精确求解的指派，与需要放松、舍入或分支定界的混合整数通信设计。}

\section{固定收益的资源指派可以精确求解}

如果每个子载波的用户收益已经固定，资源指派是一个组合结构很好的问题。令 $s_{kn}\in\{0,1\}$ 表示用户 $k$ 是否获得子载波 $n$，$c_{kn}\geq0$ 为这一指派产生的固定收益，例如给定功率下的速率。最简单的模型为
\begin{equation}
 \max_{\bm s}\quad \sum_{k,n}c_{kn}s_{kn},\qquad
 \sum_ks_{kn}\leq1\ (\forall n),\qquad s_{kn}\in\{0,1\}.
 \label{eq:simple-assignment}
\end{equation}
每个子载波独立分给收益最大的用户即可。若再要求每个用户最多分到一个子载波，即 $\sum_ns_{kn}\leq1$，问题成为二分图最大权匹配，可用指派算法精确求解。

在这个匹配模型中，约束矩阵是全幺模矩阵，即每个方形子矩阵的行列式均为 $0$ 或 $\pm1$。整数右端项使线性规划放松的极点为整数，因此把二元限制换为 $0\leq s_{kn}\leq1$ 后，总存在整数最优极点。若数值求解器返回分数解，它可能是多个最优整数解的凸组合；可以用匹配算法取回整数最优解。全幺模性是一个需要证明的矩阵性质，不能因为约束全是线性的就推断整数性。

\begin{example}[两个用户、两个资源]
固定收益矩阵及两种完整一一指派如下，收益单位可以是 bit/s/Hz：
\begin{center}
\begin{tabular}{c|cc}
\toprule
 & 资源 $1$ & 资源 $2$\\
\midrule
用户 $1$ & $8$ & $7$\\
用户 $2$ & $6$ & $2$\\
\bottomrule
\end{tabular}
\qquad
\begin{tabular}{c|c}
\toprule
指派 & 总收益\\
\midrule
$(1\to1,\ 2\to2)$ & $8+2=10$\\
$(1\to2,\ 2\to1)$ & $7+6=13$\\
\bottomrule
\end{tabular}
\end{center}
若先把最大单项 $8$ 固定给用户 $1$，剩余收益只有 $2$；联合指派得到 $13$。贪心选择最大单项没有处理不同用户对资源的竞争。这里线性规划放松也取得 $13$，可作为精确最优证书。
\end{example}

\section{联合指派与功率：时间共享放松}

当 $c_{kn}$ 本身依赖功率、干扰、天线选择或基站关联时，前面的固定收益结构不再直接适用。以没有跨子载波干扰的 OFDMA 为例，令 $e_{kn}\geq0$ 为一个调度周期内归一化的平均发射功率，$g_{kn}>0$ 为包含噪声归一化后的增益。二元 $s_{kn}$ 为 $1$ 时，该用户使用整个周期；为 $0$ 时应有 $e_{kn}=0$。常用联立约束 $e_{kn}\leq P_{kn}^{\max}s_{kn}$ 排除未指派仍发射。

将 $s_{kn}$ 放松至 $[0,1]$ 后，可以解释为使用时长的比例。使用期间的瞬时功率为 $e_{kn}/s_{kn}$，对应平均速率
\begin{equation}
 r_{kn}(s_{kn},e_{kn})=
 s_{kn}\log_2\left(1+\frac{g_{kn}e_{kn}}{s_{kn}}\right),\qquad s_{kn}>0.
 \label{eq:ofdma-perspective}
\end{equation}
这是凹函数的透视变换，对 $(s_{kn},e_{kn})$ 联合凹\cite{boyd2004}。采用 $e_{kn}\leq P_{kn}^{\max}s_{kn}$ 时，$s_{kn}=0$ 必然使 $e_{kn}=0$，闭包定义取 $r_{kn}(0,0)=0$。这样，最大化加权总速率、附上线性总功率约束与 $\sum_ks_{kn}\leq1$，便得到凸优化。

这个放松的物理意义取决于系统能否在一个调度周期内细分资源。允许时间共享时，分数 $s_{kn}$ 是可实施方案；若每个子载波必须在整个周期内只服务一个用户，它只提供最大化问题的上界。四舍五入后要重新优化功率，并检查最低速率等约束；简单舍入不保持原可行性。

\begin{example}[线性附加约束可以破坏整数性]
在式~\eqref{eq:simple-assignment} 的一一匹配版本中，再加约束 $s_{11}=s_{22}$ 和 $s_{11}+s_{22}\leq1$，则 $s_{11}=s_{22}=1/2$、$s_{12}=s_{21}=0$ 是放松的一个分数极点：两条新增约束与两个零值约束已唯一确定这四个变量。而这两个对角变量的整数可行取值只能同时为零。若只奖励对角指派，放松有正收益而整数问题只能取得零收益。新增线性约束会改变极点结构，因此原指派矩阵的整数性不能自动继承到扩展模型。
\end{example}

\section{分支定界如何形成全局证书}

分支定界通过解一系列较小的放松问题，逐步缩小最优值区间。以最大化为例，一个原问题可行方案给出下界 $L$；放松问题允许更多方案，其最优值给出上界 $U$。只有可行方案的目标才能更新 $L$，只有经过有效松弛得到的界才能用于剪枝。

在某个节点上，部分整数变量已固定。先解该节点的连续放松：若不可行，删除该节点；若其上界不超过当前 $L$，该分支不可能改进已有方案，也可删除；若得到整数可行解，则更新 $L$；否则选择某个分数变量，把它分别固定为 $0$ 和 $1$ 形成两个子节点。对一般整数变量则在其分数值的向下、向上取整处划分。这个过程反复进行，直到所有活跃节点都被处理，或全局界差满足容差。

\begin{figure}[htbp]
\centering
\begin{tikzpicture}[every node/.style={draw,rounded corners,align=center,font=\small},
 level 1/.style={sibling distance=6.5cm,level distance=1.7cm},
 edge from parent/.style={draw,-{Stealth}}]
\node {原问题的连续放松\\上界 $U$，分数变量 $s_{kn}$}
 child {node {$s_{kn}=0$\\解本节点放松、更新上界}}
 child {node {$s_{kn}=1$\\解本节点放松、更新上界}};
\end{tikzpicture}
\caption{二元变量的一次分支。子节点分别限制该变量取零或一，覆盖父节点的全部整数方案；节点放松的最优值提供最大化问题的上界。}
\label{fig:branching}
\end{figure}

对于尚未结束的搜索，全局上界取所有活跃节点上界与当前可行下界的最大值。可以报告
\begin{equation}
 L\leq R^*\leq U,\qquad
 \text{相对界差}=\frac{U-L}{\max\{1,\abs L\}},
 \label{eq:mip-gap}
\end{equation}
并写明绝对界差的单位。若已求得 $L=13$、所有未搜索节点上界至多为 $13.2$，则最优值距已有方案至多 $0.2$；这比只报告“算法收敛到 $13$”提供了更明确的信息。

混合整数\emph{凸}问题可用凸放松构造节点界。若节点中仍有非凸干扰设计，仅用 WMMSE 找到一个局部解不能作为该节点的最大化上界；它至多在满足全部约束后贡献一个可行下界。要保证全局性，节点还需半正定、凸包、区间或其他有效松弛。即使每个节点界都有效，最坏情况下仍可能搜索指数多个节点，因此小实例常用于检验大规模启发式算法的差距。

\section{如何在精确解与近似解之间选择}

固定收益的独立指派先用直接选择，带一一匹配结构时用匹配算法；加入一般耦合整数限制后，再考虑混合整数求解器。若连续放松对应可实施的时间共享，可以直接采用其解；若业务强制离散决策，就从放松中恢复可行方案并给出上界差。

整数变量与连续变量交替更新通常能迅速改善方案，但不产生通用最优证书。对于通信系统设计，先用小规模分支定界校验建模、目标尺度和启发式质量，再研究大规模实现，通常比只展示一条局部收敛曲线更有说服力。
